%! The Fundamental Theorem of Linear Algebra — Unit 3, Lesson 3.6 — Group Activity (Answer Key)
\documentclass[10pt]{article}
\usepackage{linalg-article}
\usepackage{linalg-key}   % pulls in -boxes; do NOT also load -boxes
\newcommand{\TallMath}[1]{$\displaystyle #1\rule[-1.4em]{0pt}{3.2em}$}
\newcommand{\vv}{\mathbf{v}}
\newcommand{\ww}{\mathbf{w}}
\newcommand{\xx}{\mathbf{x}}
\newcommand{\yy}{\mathbf{y}}
\newcommand{\aaa}{\mathbf{a}}
\newcommand{\svec}{\mathbf{s}}
\newcommand{\zero}{\mathbf{0}}
\newcommand{\T}{\mathsf{T}}

\begin{document}

\pageheader{Unit 3, Lesson 3.6}{Group Activity --- Answer Key}
\namepartnerperiod

\vspace{0.10in}

\begin{headlinebox}{blush}
\small\textbf{The Big Picture.} The Fundamental Theorem has two parts. \textbf{Part~1:} the four dimensions
$r,\,n-r,\,r,\,m-r$. \textbf{Part~2:} in each room the two subspaces are \emph{perpendicular} --- $N(A)\perp
C(A^{\T})$ in $\mathbb{R}^n$ and $N(A^{\T})\perp C(A)$ in $\mathbb{R}^m$. You check a right angle the same way
as always: a \emph{dot product of $0$}. Work with your partner and \emph{justify} every answer.
\end{headlinebox}

\vspace{0.08in}

\begin{tcolorbox}[colback=white, colframe=black!40, title=\textbf{Tier R --- Remediate},
    fonttitle=\bfseries, arc=1.5mm, left=3mm, right=3mm, top=2mm, bottom=2mm, breakable]
A matrix $A$ is $3\times 4$ with rank $r=2$.
\begin{enumerate}[leftmargin=1.7em, itemsep=8pt]
  \item \textbf{Two rooms.} The input room is $\mathbb{R}^{\ans{$4$}}$ and the output room is
        $\mathbb{R}^{\ans{$3$}}$.
  \item \textbf{Place each subspace} in its room (write ``input'' or ``output'') and give its dimension.
    \begin{enumerate}[label=(\alph*), leftmargin=1.8em, itemsep=4pt]
      \item Column space $C(A)$: room \ans{output ($\mathbb{R}^3$)}, $\dim=\ans{$2$}$
      \item Row space $C(A^{\T})$: room \ans{input ($\mathbb{R}^4$)}, $\dim=\ans{$2$}$
      \item Nullspace $N(A)$: room \ans{input ($\mathbb{R}^4$)}, $\dim=\ans{$2$}$
      \item Left nullspace $N(A^{\T})$: room \ans{output ($\mathbb{R}^3$)}, $\dim=\ans{$1$}$
    \end{enumerate}
  \item \textbf{The two right angles.} Name the perpendicular pair in each room.
        \\ In $\mathbb{R}^n$: \ans{row space $C(A^{\T})$} $\perp$ \ans{nullspace $N(A)$}
        \\ In $\mathbb{R}^m$: \ans{column space $C(A)$} $\perp$ \ans{left nullspace $N(A^{\T})$}
\end{enumerate}
\end{tcolorbox}

\begin{tcolorbox}[colback=white, colframe=black!40, title=\textbf{Tier A --- Approaching Proficiency},
    fonttitle=\bfseries, arc=1.5mm, left=3mm, right=3mm, top=2mm, bottom=2mm, breakable]
For each matrix (reduction supplied): give a basis for all four subspaces, then \emph{verify both right
angles} by computing dot products against the bases.
\begin{enumerate}[leftmargin=1.7em, itemsep=9pt]
  \item $A=\begin{bmatrix} 1 & 2 & 1 \\ 1 & 3 & 2 \\ 2 & 5 & 3 \end{bmatrix}
        \longrightarrow \begin{bmatrix} 1 & 0 & -1 \\ 0 & 1 & 1 \\ 0 & 0 & 0 \end{bmatrix}$
        \; (nullspace $\svec=\begin{bsmallmatrix} 1\\-1\\1 \end{bsmallmatrix}$, left null
        $\yy=\begin{bsmallmatrix} 1\\1\\-1 \end{bsmallmatrix}$).
        \ansline{$m=n=3$, $r=2$. \textbf{Bases:} $C(A^{\T})$: $(1,0,-1),(0,1,1)$; $N(A)$: $\svec=\begin{bsmallmatrix} 1\\-1\\1 \end{bsmallmatrix}$; $C(A)$: pivot columns $\aaa_1=\begin{bsmallmatrix} 1\\1\\2 \end{bsmallmatrix},\aaa_2=\begin{bsmallmatrix} 2\\3\\5 \end{bsmallmatrix}$; $N(A^{\T})$: $\yy=\begin{bsmallmatrix} 1\\1\\-1 \end{bsmallmatrix}$. \textbf{Right angle in $\mathbb{R}^3$ (input):} $(1,0,-1)\cdot\svec=1+0-1=0$, $(0,1,1)\cdot\svec=0-1+1=0$, so $N(A)\perp C(A^{\T})$. \textbf{Right angle in $\mathbb{R}^3$ (output):} $\yy\cdot\aaa_1=1+1-2=0$, $\yy\cdot\aaa_2=2+3-5=0$, so $N(A^{\T})\perp C(A)$.}
  \item $A=\begin{bmatrix} 1 & 1 & 2 & 1 \\ 1 & 2 & 3 & 1 \\ 2 & 3 & 5 & 2 \end{bmatrix}
        \longrightarrow \begin{bmatrix} 1 & 0 & 1 & 1 \\ 0 & 1 & 1 & 0 \\ 0 & 0 & 0 & 0 \end{bmatrix}$
        \; (left null $\yy=\begin{bsmallmatrix} 1\\1\\-1 \end{bsmallmatrix}$; \emph{two} nullspace
        vectors).
        \ansline{$m=3$, $n=4$, $r=2$. \textbf{Bases:} $C(A^{\T})\subseteq\mathbb{R}^4$: $(1,0,1,1),(0,1,1,0)$; $N(A)\subseteq\mathbb{R}^4$: $\svec_1=\begin{bsmallmatrix} -1\\-1\\1\\0 \end{bsmallmatrix},\svec_2=\begin{bsmallmatrix} -1\\0\\0\\1 \end{bsmallmatrix}$; $C(A)\subseteq\mathbb{R}^3$: $\aaa_1=\begin{bsmallmatrix} 1\\1\\2 \end{bsmallmatrix},\aaa_2=\begin{bsmallmatrix} 1\\2\\3 \end{bsmallmatrix}$; $N(A^{\T})\subseteq\mathbb{R}^3$: $\yy=\begin{bsmallmatrix} 1\\1\\-1 \end{bsmallmatrix}$. \textbf{Input right angle:} $(1,0,1,1)\cdot\svec_1=-1+1=0$, $(1,0,1,1)\cdot\svec_2=-1+1=0$, $(0,1,1,0)\cdot\svec_1=-1+1=0$, $(0,1,1,0)\cdot\svec_2=0$ --- so $N(A)\perp C(A^{\T})$. \textbf{Output right angle:} $\yy\cdot\aaa_1=1+1-2=0$, $\yy\cdot\aaa_2=1+2-3=0$ --- so $N(A^{\T})\perp C(A)$.}
\end{enumerate}
\end{tcolorbox}

\begin{tcolorbox}[colback=white, colframe=black!40, title=\textbf{Tier E --- Extension},
    fonttitle=\bfseries, arc=1.5mm, left=3mm, right=3mm, top=2mm, bottom=2mm, breakable]
\textbf{Why the right angles are forced.}
\begin{enumerate}[leftmargin=1.7em, itemsep=9pt]
  \item \textbf{The one-line proof.} Explain, using ``$A\xx$ is computed row by row,'' why every nullspace
        vector $\xx$ is perpendicular to every row of $A$ --- and hence to the whole row space.
        \ansline{Entry $i$ of $A\xx$ is (row$_i$ of $A$)$\,\cdot\,\xx$. If $A\xx=\zero$, every one of those dot products is $0$, so $\xx$ is perpendicular to each row. A dot product is linear, so $\xx$ is also perpendicular to every combination of the rows --- i.e.\ to the entire row space. Hence $N(A)\perp C(A^{\T})$.}
  \item \textbf{Meet only at $\zero$.} Suppose a vector $\vv$ lies in \emph{both} the row space and the
        nullspace. Then $\vv\perp\vv$, so $\vv\cdot\vv=0$. What does that force $\vv$ to be, and why? (Recall
        $\vv\cdot\vv=\|\vv\|^2$ from Lesson~1.2.)
        \ansline{$\vv\cdot\vv=\|\vv\|^2$, and this is $0$ only when $\|\vv\|=0$, i.e.\ $\vv=\zero$. So the only vector in both the row space and the nullspace is the zero vector --- the two subspaces meet \emph{only} at $\zero$. (Together with their dimensions adding to $n$, that is what ``orthogonal complements'' means.)}
  \item \textbf{The invertible case.} Let $A$ be an invertible $n\times n$ matrix ($r=n$). State the four
        dimensions, and say what Part~2 (the two right angles) reduces to here.
        \ansline{$r=n=m$: $\dim C(A)=\dim C(A^{\T})=n$, and $\dim N(A)=\dim N(A^{\T})=0$. Both nullspaces are just $\{\zero\}$, so each row/column space fills its whole room ($\mathbb{R}^n$). Part~2 becomes trivial: the only thing perpendicular to everything is $\zero$, and indeed $N(A)=N(A^{\T})=\{\zero\}$.}
\end{enumerate}
\end{tcolorbox}

\begin{teachernote}
Tier~A is the core skill: from a supplied reduction, list all four bases and then \emph{verify both right
angles} by dot products. Insist students check against \emph{every} basis vector (both reduced rows, both
nullspace vectors in problem~2), not just one --- that is the most common shortcut error. Problem~2's two
nullspace vectors make the point that ``perpendicular to a subspace'' means perpendicular to a whole basis.
Tier~E~1 is the must-do proof (row-by-row $\Rightarrow$ perpendicular); Tier~E~2 is the elegant reason the
pair meets only at $\zero$ ($\vv\cdot\vv=0\Rightarrow\vv=\zero$), which is worth sharing with the whole class
if time allows.
\end{teachernote}

\end{document}
